\Section{Skills, Hooks, and MCP Integration}
\label{sec:skills}

The agentic side of the workflow is delivered as a Claude Code plugin, packaging skills (markdown instructions), edit-time hooks, and an MCP server.
The plugin is generated from Tera templates so that shared content (such as the verification sub-workflow) is included by reference rather than duplicated.

\Paragraph{Skills}
The inference skill registers the workflow of Sec.~\ref{sec:example} as a fixed task sequence --- synthesize loop invariants, run \WP, simplify, full-scope verification, fix logical errors, resolve timeouts --- and tags every emitted condition with an \texttt{[inferred]} marker so that re-runs are idempotent and never disturb user-written specs.
Beyond the task script, the skills concentrate the domain guidance the agent would otherwise have to rediscover.
\emph{Loop invariants:} when to introduce recursive helper functions for iterative quantities, and when to prefer state-level invariants on resources --- assumed by the prover at every call site, including each loop iteration --- over per-loop inductive arguments.
\emph{Simplification:} discard vacuous conditions; replace unbounded quantifiers, a frequent source of SMT timeouts, with bounded or closed-form equivalents; and require trigger annotations on the surviving ones.
\emph{Timeout resolution:} a hierarchy of escalations --- strengthen state-level invariants, then introduce spec helpers, then lemmas, then explicit proof scripts --- with a strict prohibition on weakening abort or post conditions to silence a timeout, since this trades a real property for the appearance of progress.
The simplification and timeout-resolution material is shared, via template inclusion, with a standalone verification skill, so the same guidance applies whether the spec was inferred or written by hand.

\Paragraph{Hooks}
Event hooks in \texttt{hooks.json} couple the plugin to the agent's life cycle.
The most important hook is \texttt{PostToolUse}, triggered after every \texttt{Edit} or \texttt{Write} on a \texttt{.move} file: it parses the modified file, runs a set of offline AST checks (e.g.\ \texttt{old()} used outside \texttt{ensures}, references or dereferences inside spec expressions, deprecated Move~1 syntax such as \texttt{borrow\_global<\ldots>}), and auto-formats.
When any check fails, the hook exits non-zero and the diagnostic is surfaced back into the agent's transcript, where it functions as immediate feedback that the next turn must address --- the agent is, in effect, given a fast linter loop without the round-trip cost of a full prover invocation.

\Paragraph{MCP integration}
The plugin also ships a stdio MCP server that exposes the Move toolchain to \CC as a small set of tools: compile-status checking, package manifest inspection, structural queries (dependency graph, module summary, call graph, function usage), invoking the prover with a per-function filter, a per-VC timeout, and an exclude list, and running the \WP-based spec inference of Sec.~\ref{sec:wp} with a function- or module-level filter.
The server holds a long-lived cache for the active package, invalidated by an OS-native file watcher on source changes.
This makes the cost of a single tool call negligible relative to a fresh \texttt{cargo build}, which is what makes the iterative \emph{infer / verify / refine} loop tractable at the granularity inference needs.

%%% Local Variables:
%%% mode: latex
%%% TeX-master: "main"
%%% End:
